% 05-methodes-numeriques.tex — État de l'art 4/5 : méthodes numériques pour
% la simulation physique interactive (≈ 4 pages).

\chapter{Méthodes numériques pour la simulation interactive}
\label{chap:methodes}

\section{Intégrer les équations du mouvement}

\begin{arediger}{1,5 page}
Euler explicite, Euler semi-implicite, Verlet \parencite{verlet1967},
Runge--Kutta d'ordre~4 \parencite{butcher2016} ; précision, stabilité,
conservation de l'énergie \parencite{hairer2006}. Justifier RK4 pour le
mouvement le long du rail (une seule équation \(\ddot{s} = -g \sin(\varphi)\))
et les formes exactes partout où elles existent (parabole, énergie).
\end{arediger}

\section{Dynamique basée sur les positions}

\begin{arediger}{1,5 page}
Principe de la dynamique basée sur les positions \parencite{muller2007pbd,
bender2014survey}, sa version étendue XPBD \parencite{macklin2016xpbd} et
l'approche à petits pas \parencite{macklin2019smallsteps} ; précurseur dans le
jeu vidéo \parencite{jakobsen2001}. Cordes, cheveux et tiges : projection
\enquote{follow-the-leader} \parencite{muller2012ftl} et solveur direct pour
chaînes \parencite{deul2018direct}. Le chapitre~\ref{chap:jeu} applique ces
méthodes.
\end{arediger}

\section{Fluides particulaires}

\begin{arediger}{1 page}
SPH \parencite{monaghan1992sph} et fluides basés sur les positions
\parencite{macklin2013pbf}. Pourquoi nous nous limitons à dix particules sans
pression ni viscosité : on veut voir \emph{quand} l'eau sort du seau, pas
comment elle s'écoule.
\end{arediger}
